% !TEX root = ../main.tex
%
% Section 4 -- Model Architecture
% Requires (in preamble of main.tex): amsmath, amssymb, booktabs, array
%   \usepackage{amsmath}
%   \usepackage{amssymb}
%   \usepackage{booktabs}
%   \usepackage{array}

\section{Methods}
\subsection{Model Architecture}

Table~\ref{tab:model-overview} summarizes the branch composition of all six models
compared in this study. Each model is defined by which of three modality-specific
branches it contains, and, for WSI-containing models, by which pooling mechanism is used to aggregate patch-level information into a slide-level representation. The three different branches are defined as WSI (Whole-Slide Image) Branch, Clinical Branch, and RNA-seq Branch by what data are used as branches' input.

\begin{table}[htbp]
  \centering
  \begin{tabular}{@{}l l l l@{}}
    \toprule
    \textbf{Model} & \textbf{WSI Branch} & \textbf{Clinical Branch} & \textbf{RNA-seq Branch} \\
    \midrule
    M1 & \checkmark (ABMIL)                  & ---                        & --- \\
    M2 & \checkmark (ABMIL)                  & \checkmark (age/sex MLP)   & --- \\
    M3 & \checkmark (RNA-seq co-attention)   & ---                        & \checkmark\\
    M4 & \checkmark (RNA-seq co-attention)   & \checkmark& \checkmark\\
    M6 & ---                                 & ---                        & \checkmark\\
    M7 & ---                                 & \checkmark& \checkmark\\
    \bottomrule
  \end{tabular}
 \caption{Model overview by branch composition. WSI branch subtype is either
    \textit{ABMIL} (gated-attention MIL pooling into a single RNA-agnostic
    slide-level vector; M1/M2) or \textit{RNA-seq co-attention} (multi-component
    pooling followed by RNA-guided cross-attention over the pooled views; M3/M4).}
    \label{tab:model-overview}

\end{table}

M1--M4 share an identical patch encoding backbone and
differ only in modality-specific pooling and fusion.
A shared set of regularization and auxiliary modules is applied on top of this
backbone, and all models are trained with the same
risk prediction head and survival objective. M6 and M7
are WSI-free baselines used to isolate the incremental contribution of the WSI branch.

\subsection{Shared Patch Encoding Backbone}
\label{subsec:backbone}

All WSI branch-based models (M1--M4) share an identical patch-level encoding
backbone. Each whole-slide image is represented as a bag of tissue patches, and each
patch $x_{ij}$ (patch $j$ of patient $i$) is passed through a frozen convolutional
encoder, ResNet50 pretrained with SwAV self-supervision on histopathology
\cite{lunit-swav}:

\begin{align}
  z_{ij} &= f_{\mathrm{CNN}}(x_{ij}) \qquad \text{(frozen)} \\
  h_{ij} &= f_{\mathrm{proj}}(z_{ij}) \qquad \text{(trainable linear projection to \texttt{embed\_dim})}
\end{align}

The projected patch embeddings, together with a learned embedding of their
normalized spatial coordinates $\mathbf{c}_{ij}$ within the slide, are passed
through a Transformer encoder with Nystr\"om-approximated self-attention, allowing
patches to exchange spatial context before pooling:

\begin{equation}
  \{t_{ij}\}_{j=1}^{N_i} = \mathrm{ViT}\big(\{h_{ij}\}_{j=1}^{N_i}, \{\mathbf{c}_{ij}\}_{j=1}^{N_i}\big)
\end{equation}

where $N_i$ is the number of patches sampled from patient $i$'s slide(s). This
backbone (frozen CNN $\rightarrow$ trainable projection $\rightarrow$
self-attention) is identical across M1--M4. The models differ only in how the
resulting patch tokens $\{t_{ij}\}$ are pooled into a slide-level representation and
how additional modalities are fused, described next.

\subsection{Modality-Specific Pooling and Fusion}
\label{subsec:pooling-fusion}

\subsubsection{M1 / M2: Gated-Attention MIL (ABMIL)}

M1 (WSI branch-only) and M2 (WSI branch + clinical) aggregate patch tokens with
gated-attention multiple instance learning \cite{ilse2018abmil}:

\begin{align}
  \mathrm{gate}_{ij} &= \tanh(W_v t_{ij}) \odot \sigma(W_u t_{ij}) \\
  \alpha_{ij} &= \mathrm{softmax}_j\big(w^\top \mathrm{gate}_{ij}\big) \\
  z_{\mathrm{wsi}} &= \sum_j \alpha_{ij}\, t_{ij}
\end{align}

M1 predicts risk directly from $z_{\mathrm{wsi}}$. M2 additionally embeds age and
sex through a small clinical encoder, $g_i = f_{\mathrm{clinical}}(\mathrm{age}_i,
\mathrm{sex}_i)$, and concatenates it with the WSI branch representation before the risk head: $q_i = [z_{\mathrm{wsi}}, g_i]$.

\subsubsection{M3 / M4: Multi-Component Pooling with RNA-Guided Co-Attention}

M3 (WSI + RNA-seq) and M4 (WSI + RNA-seq + clinical) replace ABMIL with
a two-stage pooling mechanism designed to avoid collapsing the slide into a single
vector \emph{before} RNA information is introduced. First, \emph{multi-component
pooling} summarizes the patch tokens into four parallel statistical views rather
than one:

\begin{equation}
  \mathrm{components}_i =
  \begin{bmatrix}
    \displaystyle\operatorname*{mean}_j(t_{ij}) \\[4pt]
    \displaystyle\operatorname*{std}_j(t_{ij}) \\[4pt]
    \displaystyle\sum_j \alpha_{ij}\, t_{ij} \;\text{(ABMIL)} \\[4pt]
    \displaystyle\operatorname*{mean}\big(t_{ij} \mid \alpha_{ij} \in \text{top } 10\%\big)
  \end{bmatrix}
  \in \mathbb{R}^{4 \times D}
\end{equation}

The RNA-seq profile $m_i$ (the literature-guided 1,500-gene z-score vector) is encoded independently,
$z_{\mathrm{rna}} = f_{\mathrm{RNA}}(m_i)$, and then acts as the \textbf{query} in a
multi-head cross-attention layer over the four components (as key/value), letting
the transcriptomic profile determine which morphological view is most relevant for
that patient:

\begin{equation}
  z_{\mathrm{wsi}} = \mathrm{CoAttention}\big(\text{query}=z_{\mathrm{rna}},\ \text{key/value}=\mathrm{components}_i\big)
\end{equation}

M3 (no clinical branch) predicts risk from $q_i = [z_{\mathrm{wsi}}, z_{\mathrm{rna}}]$;
M4 additionally concatenates the clinical embedding:
$q_i = [z_{\mathrm{wsi}}, z_{\mathrm{clinical}}, z_{\mathrm{rna}}]$. This RNA-guided
co-attention is the key architectural distinction from M1/M2: instead of RNA being
appended after a fixed WSI branch summary, it directly conditions \emph{which}
WSI view is pooled.

\subsection{Regularization and Auxiliary Modules}
\label{subsec:regularization}

Four additional modules are applied on top of the architectures above as the final
training recipe.

\begin{description}
  \item[Patch sub-sampling.] At each training epoch, a random fraction
    $\kappa$ of a slide's patches is sampled before the forward pass; validation,
    internal test, and external test always use the full patch set. This
    regularizes against overfitting on the limited patient cohort.

  \item[Attention dispersion.] A parameter-free scalar summarizing how
    spatially concentrated or dispersed the attention-weighted patches are,
    computed as the attention-weighted standard deviation of patch coordinates:
    \begin{equation}
      \mathrm{disp}_i = \operatorname*{std}_j\big(\mathbf{c}_{ij};\ \text{weights}=\alpha_{ij}\big)
    \end{equation}
    This scalar scaled by a single learnable factor is appended as an additional
    view to the risk-head input $q_i$ for all four WSI branch models (M1--M4).

  \item[RNA-prediction auxiliary task (M3/M4 only).] A lightweight MLP head
    predicts the patient's 1,500-gene expression vector from the RNA-\emph{free}
    mean-pooled patch representation, $\bar{t}_i = \operatorname*{mean}_j(t_{ij})$,
    taken \emph{before} any RNA-conditioned pooling to avoid circularity/leakage,
    following the HE2RNA formulation \cite{schmauch2020he2rna}. This auxiliary
    regression is trained jointly with the main survival objective, weighted by
    $\lambda_{\mathrm{aux}}$. It applies only to M3/M4,
    since M1/M2 have no RNA encoder to supervise against.

  \item[Tile augmentation.] Training patches are augmented with random horizontal/vertical flip, color jitter, Gaussian blur at every epoch; validation, internal test, and external test are always evaluated
    on unaugmented tiles.
\end{description}

\subsection{Risk Prediction Head and Training Objective}
\label{subsec:risk-head}

All models output a single patient-level risk score through a shared head
architecture (LayerNorm $\rightarrow$ Linear $\rightarrow$ scalar):

\begin{equation}
  r_i = f_{\mathrm{head}}(q_i)
\end{equation}

where $q_i$ is the concatenated representation described above. Models are trained by
minimizing the Cox partial-likelihood negative log-likelihood, with risk sets
constructed within each training batch:

\begin{equation}
  \mathcal{L}_{\mathrm{cox}} = -\sum_{i:\, E_i = 1} \left[ r_i - \log \sum_{j:\, T_j \geq T_i} \exp(r_j) \right]
\end{equation}

where $T_i$ and $E_i$ denote patient $i$'s observed time and event indicator,
respectively. For M3/M4, the auxiliary RNA-prediction loss (mean squared error
between predicted and observed z-scored gene expression, $\mathcal{L}_{\mathrm{aux}}$)
is added to the training objective:

\begin{equation}
  \mathcal{L} = \mathcal{L}_{\mathrm{cox}} + \lambda_{\mathrm{aux}} \cdot \mathcal{L}_{\mathrm{aux}}
\end{equation}

\subsection{WSI-Free Baselines (M5, M6, M7)}
\label{subsec:wsi-free}

To isolate the contribution of the WSI branch, three WSI-free baselines are
trained: \textbf{M5} (clinical only, age/sex embedded through the same clinical
encoder as M2/M4 $\rightarrow$ risk head, no RNA-seq or WSI input), \textbf{M6}
(RNA-seq only, on the literature-guided 1,500-gene input, MLP encoder
$\rightarrow$ risk head), and \textbf{M7} (RNA-seq + clinical, concatenated
$\rightarrow$ risk head). None of the three contains a CNN/ViT backbone or the modules described above, since these are
WSI-branch-specific. Together they serve as non-WSI reference points against
which the incremental contribution of M1--M4's image branch is evaluated under
external (cross-institution) validation.